\begin{figure}[h]
  \centering
  \begin{adjustbox}{max width=\linewidth}
    \begin{tikzpicture}
      \foreach \value [count=\i from 0] in {89, 2, 3} {
        \pgfmathtruncatemacro{\at}{4 * \i}
        \ifnum\i=0
          \node[memhit] (a\i) at (\i * 1.0, 0) {\value};
        \else
          \node[memcell] (a\i) at (\i * 1.0, 0) {\value};
        \fi
        \node[memnote, above=0.5mm of a\i] {+\at};
      }
      \node[memcell, fill=black!8, right=0mm of a2] (pad) {};
      \node[memnote, above=0.5mm of pad] {+12};
      \node[memnote, right=1mm of pad] (padn) {padding};
      \node[memaddr, minimum width=20mm] (b) at (0.5, -1.4) {};
      \node[memnote, anchor=south] at ([xshift=5mm, yshift=0.5mm]b.north west) {+16};
      \node[cflab, left=3mm of a0] (an) {\token{a}};
      \node[cflab, left=3mm of b] (bn) {\token{b}};
      \node[memframe, fit=(an)(padn)(bn)(b), inner ysep=4mm] (main) {};
      \node[memtitle] at ([xshift=2mm]main.north west) {main};
      \draw[mempoint] ([xshift=-5mm]b.center) -- (a0.south);

      \node[memregion] at ($(main.north) + (0, 5mm)$) {stack};
    \end{tikzpicture}
  \end{adjustbox}
  \caption{\label{fig:memory:ref_variable} The frame of the listing above, after
    \token{a[0] = 89;}. The array \token{a} holds its 12 bytes of elements, and
    the reference \token{b} holds the address of \token{a}, aligned on 8 bytes:
    reading \token{b[0]} reads \token{a[0]}.}
\end{figure}
